
\documentclass{article}
\usepackage{colortbl}
\usepackage{makecell}
\usepackage{multirow}
\usepackage{supertabular}

\begin{document}

\newcounter{utterance}

\centering \large Interaction Transcript for game `cladder', experiment `full\_v1.5\_default', episode 1084 with refineonly{-}Qwen3.5{-}4B.
\vspace{24pt}

{ \footnotesize  \setcounter{utterance}{1}
\setlength{\tabcolsep}{0pt}
\begin{supertabular}{c@{$\;$}|p{.15\linewidth}@{}p{.15\linewidth}p{.15\linewidth}p{.15\linewidth}p{.15\linewidth}p{.15\linewidth}}
    \# & \multicolumn{2}{c}{Player} && \multicolumn{2}{c}{Game Master} \\
    \hline

    \theutterance \stepcounter{utterance}  
    & & & \multicolumn{4}{p{0.6\linewidth}}{
        \cellcolor[rgb]{0.9,0.9,0.9}{
            \makecell[{{p{\linewidth}}}]{
                \texttt{\tiny{[P1$\langle$GM]}}
                \texttt{You are an expert in causal inference. The following question is not a typical commonsense query, but rather a meticulously designed question created by a professor specializing in causal inference, intended to assess the students' mastery of the course content.} \\
\texttt{Imagine a self{-}contained, hypothetical world with only the following conditions, and without any unmentioned factors or causal relationships: Age has a direct effect on hospital costs and recovery. Hospital costs has a direct effect on recovery. The overall probability of high hospital bill is 51\%. The probability of low hospital bill and recovery is 37\%. The probability of high hospital bill and recovery is 17\%. Is the chance of recovery larger when observing high hospital bill?} \\
\texttt{Start your answer with 'Yes' or 'No', followed by additional reasoning or evidence to support your explanation.} \\
            }
        }
    }
     \\ \\

    \theutterance \stepcounter{utterance}  
    & \multicolumn{4}{p{0.6\linewidth}}{
        \cellcolor[rgb]{0.561,0.737,0.561}{
            \makecell[{{p{\linewidth}}}]{
                \texttt{\tiny{[P1$\rangle$GM]}}
                \texttt{No} \\
\\ 
\texttt{**Reasoning:**} \\
\\ 
\texttt{To determine whether the chance of recovery is larger when observing a high hospital bill, we must calculate the conditional probability of recovery given a high bill, denoted as \$P(\textbackslash{}text\{Recovery\} \textbackslash{}mid \textbackslash{}text\{High Bill\})\$.} \\
\\ 
\texttt{We are given the following joint probabilities:} \\
\texttt{1.  \$P(\textbackslash{}text\{High Bill\}) = 0.51\$} \\
\texttt{2.  \$P(\textbackslash{}text\{Low Bill, Recovery\}) = 0.37\$} \\
\texttt{3.  \$P(\textbackslash{}text\{High Bill, Recovery\}) = 0.17\$} \\
\\ 
\texttt{First, we determine the total probability of recovery, \$P(\textbackslash{}text\{Recovery\})\$. Since the bill outcomes are mutually exclusive (Low vs. High), the total probability of recovery is the sum of recovery occurring with a low bill and recovery occurring with a high bill:} \\
\texttt{\$\$P(\textbackslash{}text\{Recovery\}) = P(\textbackslash{}text\{Low Bill, Recovery\}) + P(\textbackslash{}text\{High Bill, Recovery\})\$\$} \\
\texttt{\$\$P(\textbackslash{}text\{Recovery\}) = 0.37 + 0.17 = 0.54\$\$} \\
\\ 
\texttt{Next, we calculate the conditional probability of recovery given a high bill using the definition of conditional probability:} \\
\texttt{\$\$P(\textbackslash{}text\{Recovery\} \textbackslash{}mid \textbackslash{}text\{High Bill\}) = \textbackslash{}frac\{P(\textbackslash{}text\{High Bill, Recovery\})\}\{P(\textbackslash{}text\{High Bill\})\}\$\$} \\
\texttt{\$\$P(\textbackslash{}text\{Recovery\} \textbackslash{}mid \textbackslash{}text\{High Bill\}) = \textbackslash{}frac\{0.17\}\{0.51\} \textbackslash{}approx 0.333\$\$} \\
\\ 
\texttt{For comparison, let us calculate the probability of recovery given a low bill:} \\
\texttt{\$\$P(\textbackslash{}text\{Recovery\} \textbackslash{}mid \textbackslash{}text\{Low Bill\}) = \textbackslash{}frac\{P(\textbackslash{}text\{Low Bill, Recovery\})\}\{P(\textbackslash{}text\{Low Bill\})\}\$\$} \\
\texttt{Since \$P(\textbackslash{}text\{Low Bill\}) = 1 {-} P(\textbackslash{}text\{High Bill\}) = 1 {-} 0.51 = 0.49\$:} \\
\texttt{\$\$P(\textbackslash{}text\{Recovery\} \textbackslash{}mid \textbackslash{}text\{Low Bill\}) = \textbackslash{}frac\{0.37\}\{0.49\} \textbackslash{}approx 0.755\$\$} \\
\\ 
\texttt{**Conclusion:**} \\
\texttt{The probability of recovery given a high bill (\$\textbackslash{}approx 33.3\textbackslash{}\%\$) is significantly **smaller** than the probability of recovery given a low bill (\$\textbackslash{}approx 75.5\textbackslash{}\%\$).} \\
\\ 
\texttt{Therefore, the answer to the question "Is the chance of recovery larger when observing high hospital bill?" is **No**. In fact, observing a high bill indicates a lower likelihood of recovery compared to observing a low bill. This counter{-}intuitive result arises because "High Bill" and "Recovery" are negatively correlated in the joint distribution (both occur together less frequently} \\
            }
        }
    }
    & & \\ \\

    \theutterance \stepcounter{utterance}  
    & & & \multicolumn{2}{p{0.3\linewidth}}{
        \cellcolor[rgb]{0.9,0.9,0.9}{
            \makecell[{{p{\linewidth}}}]{
                \texttt{\tiny{[GM$|$GM]}}
                \texttt{No} \\
            }
        }
    }
    & & \\ \\

    \theutterance \stepcounter{utterance}  
    & & & \multicolumn{2}{p{0.3\linewidth}}{
        \cellcolor[rgb]{0.9,0.9,0.9}{
            \makecell[{{p{\linewidth}}}]{
                \texttt{\tiny{[GM$|$GM]}}
                \texttt{no} \\
            }
        }
    }
    & & \\ \\

    \theutterance \stepcounter{utterance}  
    & & & \multicolumn{2}{p{0.3\linewidth}}{
        \cellcolor[rgb]{0.9,0.9,0.9}{
            \makecell[{{p{\linewidth}}}]{
                \texttt{\tiny{[GM$|$GM]}}
                \texttt{game\_result = WIN} \\
            }
        }
    }
    & & \\ \\

\end{supertabular}
}

\end{document}
